\subsection*{4.4.2 Handling Well-Behaved Buckets}
We will use our tail bounds for subset-wise symmetric polynomials to prove the following lemma.
\begin{authorclaim}{Lemma17}
With probability at least $1-1/(2n)$ over $z\sim Z$, $f_{\mathcal B}$ can be written as $f_{\mathcal B}=f'_{\mathcal B}\oplus f''_{\mathcal B}$, where $f'_{\mathcal B}$ and $f''_{\mathcal B}$ are over disjoint set of variables, and for every $g\in\{f'_{\mathcal B},f''_{\mathcal B}\}$ it holds that $|\mathbb E[\widetilde g_z(Y)]-\mathbb E[\mathfrak g]|\le1/n$.
\end{authorclaim}
The proof of Lemma17 will follow immediately from Claim20 and Lemma21. Toward proving the above lemma, let us set some preliminaries.
\begin{authorclaim}{Claim18}
If $q\in\mathcal B$ then $c^q\in[C\log\log n,C\log n]$ and $m_q\le3^{c^q}$.
\end{authorclaim}
\begin{proof}
The upper bound on $c^q$ follows immediately from the assumption in Lemma15 that every term has width at most $C\log n$. Also, $m_q>\log^{2C}n$ since $q\notin\mathcal A$. Since we are at CaseII, $m_q>\log^{2C}n$ implies that $m_q\le3^{c^q}$. From the fact that $\log^{2C}n<3^{c^q}$ we get $c^q>\log_3(\log^{2C}n)>C\log\log n$.
\end{proof}
Recall that a term $f_i$ is good under $z$ if the variables read by $f_i$ intersects with $z$ in at least $1-\beta$ fraction.
\begin{authorclaim}{Claim19}
For a fixed $z\in\{0,1\}^n$, let $X_z\subseteq[m]$ be the set of terms in $f_{\mathcal B}$ that are not good under $z$. Then, with probability at least $1-1/(2n)$ over $z\sim Z$, $|X_z|\le(3c/(c-1))\log n$.
\end{authorclaim}
\begin{proof}
The proof is very similar to Claim13. Fix a bucket $q\in\mathcal B$, set $k_q=3\log n/c^q$ and observe that $k\ge k_q$. Let $S\subseteq B_q$ be some subset of cardinally $k_q$. We first bound the probability $p$ that every term $f_i$ for $i\in S$ is bad under $z\sim Z$.
For a truly random $z\sim\operatorname{Bernoulli}(1-2^{-C})^{\otimes n}$, the above probability is bounded by
\[
\prod_{i\in S}\binom{w_i}{\beta w_i}2^{-\beta Cw_i}
\le\prod_{i\in S}2^{w_i}2^{-5w_i}
\le(2^{-4c^{q-1}})^{k_q}
\le2^{-3k_qc^q}\le n^{-9}.
\]
For $z\sim Z$, we get that $p\le n^{-9}+\gamma\le2n^{-9}$. Thus, with probability at most $\binom{m_q}{k_q}p$ over $z\sim Z$ there exists a set of $k_q$ terms in $B_q$ whose all terms are bad under $z$. By using Claim18, we get
\[
\binom{m_q}{k_q}p\le m_q^{k_q}\cdot2^{-9\log n+1}
\le3^{k_qc^q+\log_32\cdot(-9\log n+1)}
\le3^{-9\log n/4}\le n^{-3}.
\]
Moreover, with probability at most $|\mathcal B|n^{-3}\le n^{-2}$ over $z\sim Z$ there exists a $q\in\mathcal B$ and a set of $k_q$ terms in $B_q$ whose all terms are bad under $z$. Taking the contrapositive, we infer that with probability at least $1-n^{-2}\ge1-1/(2n)$ over $z\sim Z$, we have at most
\[
\sum_{q\in\mathcal B}k_q\le\sum_{q=1}^Q\frac{3\log n}{c^q}\le\frac{3c}{c-1}\log n
\]
terms that are bad for $z$.
\end{proof}
From here onwards, we fix a $z$ satisfying $|X_z|\le(3c/(c-1))\log n$. Write
\[
f_{\mathcal B}=\bigoplus_{i\in\mathcal C\setminus X_z}f_i\oplus\bigoplus_{i\in X_z}f_i
\mathrel{\mathop:}=f'_{\mathcal B}\oplus f''_{\mathcal B},
\]
where $\mathcal C=\bigcup_{q\in\mathcal B}B_q\subseteq[m]$ is the set of all terms that belong to $\mathcal B$'s buckets. Simply put, we divide $f_{\mathcal B}$ to the parity of exceptional terms $f''_{\mathcal B}$ and non-exceptional terms $f'_{\mathcal B}$ for whom we will refer to as good terms. We stress that both $f'_{\mathcal B}$ and $f''_{\mathcal B}$ depend on $z$.
\begin{authorclaim}{Claim20 (Exceptional terms)}
$|\mathbb E[(\widetilde{f''_{\mathcal B}})_z(Y)]-\mathbb E[\mathfrak f''_{\mathcal B}]|\le1/n$.
\end{authorclaim}
\begin{proof}
For brevity, let $g=f''_{\mathcal B}$. For a fixed $w\in\{0,1\}^n$, let $g^w(x)=\mathfrak g(x+w)$. The proof will follow from bounding the spectral norm of $g^w$. Indeed, $g^w$ is a multiplication of at most $(3c/(c-1))\log n$ terms, each of which has spectral norm at most3. By sub-multiplicativity,
\[
L_1(g^w)\le3^{(3c/(c-1))\log n}\le n^{5c/(c-1)}.
\]
Now, $\delta\le\frac12n^{-5c/(c-1)-1}$, so by Claim5 we get that $|\mathbb E[g^w(Y)]-\mathbb E[g^w]|\le1/n$ for every $w\in\{0,1\}^n$.
Fooling $g^w$ is sufficient to fool $\widetilde g_z$. To see this, note that
\begin{align*}
|\mathbb E[\widetilde g_z(Y)]-\mathbb E[\mathfrak g]|
&=|\mathbb E[\mathfrak g(Y+z\wedge U)]-\mathbb E[\mathfrak g(U+z\wedge U')]|\\
&=\left|\mathbb E_{w\sim U}[\mathbb E[g^w(Y)]-\mathbb E[g^w]]\right|\le1/n,
\end{align*}
where $U'$ is an independent copy of $U$.
\end{proof}
Next, we prove:
\begin{authorclaim}{Lemma21 (Good terms)}
$|\mathbb E[(\widetilde{f'_{\mathcal B}})_z(Y)]-\mathbb E[\mathfrak f'_{\mathcal B}]|\le1/n$.
\end{authorclaim}
\begin{proof}
For brevity, let $g=f'_{\mathcal B}$ and recall that its set of terms is given by $\mathcal C\setminus X_z$. Shifting the bias function $\widetilde g=\widetilde g_z$ to mean zero, recall that we define $\check g(x)=\widetilde g(x)/\mathbb E[\widetilde g]-1$.
Thus, we can write
\begin{align*}
\widetilde g
&=\mathbb E[\mathfrak g]\prod_{i\in\mathcal C\setminus X_z}(1+\widehat g_i)
=\mathbb E[\mathfrak g]\sum_{I\subseteq\mathcal C\setminus X_z}\prod_{i\in I}\widehat g_i\\
&=\mathbb E[\mathfrak g]\sum_{\vec k\in\mathbb N^Q}
\sum_{\substack{I\subseteq\mathcal C\setminus X_z\\K(I)=\vec k}}\prod_{i\in I}\widehat g_i,
\end{align*}
where by $K(I)=\vec k$ we mean that for every $q\in[Q]$, there are $\vec k[q]$ terms in $I$ that belong to the $q$-th bucket, i.e., $|I\cap B_q|=\vec k[q]$. For simplicity, we reorder the terms of $g$ and write $g=\bigoplus_{i\in[m']}g_i$ for $m'=|\mathcal C\setminus X_z|$, and for $q\in[Q]$, $B_q\subseteq[m']$ is the set of terms in $g$ that belong to the $q$-th bucket. We abbreviate $\vec g=(\widehat g_1,\ldots,\widehat g_{m'})$, and write
\[
S_{\vec k}(\vec g)=\sum_{\substack{I\subseteq[m']\\K(I)=\vec k}}\prod_{i\in I}\widehat g_i.
\]
Under these notations, $\widetilde g=\mathbb E[\mathfrak g]\sum_{\vec k\in\mathbb N^Q}S_{\vec k}(\vec g)$.
Let $I^g(x)$ be the Boolean-valued function which is1 if and only if
\[
\sum_{\substack{\vec k\in\mathbb N^Q\\\|\vec k\|_{(c)}>A}}|S_{\vec k}(\vec g(x))|\le2^{-A/1024},
\]
where $A=\bar C\log n$ and $\|\vec k\|_{(c)}=\sum_{q=1}^Qc^q\vec k[q]$. Section4.4.4 will be devoted to showing that $\mathbb E[I^g(Y)]$ is very close to1. Namely,
\begin{authorclaim}{Lemma22}
The following two inequalities hold.
\begin{enumerate}
\item $\mathbb E[I^g(Y)]\ge1-e^{-c_IA}$ for $c_I=\ln2/2^{23}$.
\item $\mathbb E[S_{\vec k}^2(\vec g(Y))]\le2^{-\|\vec k\|_{(c)}/8}$.
\end{enumerate}
\end{authorclaim}
For now, let us take Lemma22 as given and continue with the proof of Lemma21. We proceed by writing
\begin{equation}
|\mathbb E[\widetilde g(Y)]-\mathbb E[\widetilde g]|
\le|\mathbb E[\widetilde g(Y)\mid I^g(Y)=1]-\mathbb E[\widetilde g]|
+2\Pr[I^g(Y)=0].\tag{8}
\end{equation}
By Lemma22, we have that $\Pr[I^g(Y)=0]\le e^{-c_IA}$. Next, observe that
\[
|\mathbb E[\widetilde g(Y)\mid I^g(Y)=1]-\mathbb E[\widetilde g]|
=\left|\mathbb E[\mathfrak g]\sum_{\substack{\vec k\in\mathbb N^Q\\\|\vec k\|_{(c)}>0}}
\mathbb E[S_{\vec k}(\vec g(Y))\mid I^g(Y)=1]\right|,
\]
and set
\[
\Delta=\left|\sum_{\substack{\vec k\in\mathbb N^Q\\\|\vec k\|_{(c)}>0}}
\mathbb E[S_{\vec k}(\vec g(Y))\mid I^g(Y)=1]\right|,
\]
so Equation(8) gives us
\begin{equation}
|\mathbb E[\widetilde g(Y)]-\mathbb E[\widetilde g]|\le\Delta+2e^{-c_IA}.\tag{9}
\end{equation}
We bound $\Delta$ as follows.
\begin{align*}
\Delta&\le\left|\sum_{\substack{\vec k\in\mathbb N^Q\\0<\|\vec k\|_{(c)}\le A}}
\mathbb E[S_{\vec k}(\vec g(Y))\mid I^g(Y)=1]\right|\\
&\quad+\max_{\substack{y\in\{0,1\}^n\\I^g(y)=1}}
\sum_{\substack{\vec k\in\mathbb N^Q\\\|\vec k\|_{(c)}>A}}|S_{\vec k}(\vec g(y))|.
\end{align*}
By definition, the second term is at most $2^{-A/1024}$. The first term, call it $\Delta_1$, can be split into two terms as follows.
\begin{align}
\Delta_1
&=\frac1{\Pr[I^g(Y)=1]}\left|\sum_{\substack{\vec k\in\mathbb N^Q\\0<\|\vec k\|_{(c)}\le A}}\mathbb E[S_{\vec k}(\vec g(Y))\cdot I^g(Y)]\right|\notag\\
&\le2\left|\sum_{\substack{\vec k\in\mathbb N^Q\\0<\|\vec k\|_{(c)}\le A}}\mathbb E[S_{\vec k}(\vec g(Y))\cdot I^g(Y)]\right|\tag{10}\\
&\le2\left|\sum_{\substack{\vec k\in\mathbb N^Q\\0<\|\vec k\|_{(c)}\le A}}\mathbb E[S_{\vec k}(\vec g(Y))]\right|
+2\left|\sum_{\substack{\vec k\in\mathbb N^Q\\0<\|\vec k\|_{(c)}\le A}}\mathbb E[S_{\vec k}(\vec g(Y))\cdot(1-I^g(Y))]\right|\notag\\
&\le2\left|\sum_{\substack{\vec k\in\mathbb N^Q\\0<\|\vec k\|_{(c)}\le A}}\mathbb E[S_{\vec k}(\vec g(Y))]\right|
+2\sqrt{\mathbb E[1-I^g(Y)]}\sum_{\substack{\vec k\in\mathbb N^Q\\0<\|\vec k\|_{(c)}\le A}}\sqrt{\mathbb E[S_{\vec k}^2(\vec g(Y))]}.\tag{11}
\end{align}
where the last inequality follows from the triangle inequality followed by Cauchy-Schwarz.
By Lemma22, the second term of Equation(11), $\Delta_{1,2}$, is at most
\begin{align*}
\Delta_{1,2}
&\le2e^{-c_IA}\sum_{\substack{\vec k\in\mathbb N^Q\\0<\|\vec k\|_{(c)}\le A}}\sqrt{2^{-\|\vec k\|_{(c)}/8}}\\
&\le2e^{-c_IA}\sum_{w=1}^{A-1}|\{\vec k\in\mathbb N^Q:w<\|\vec k\|_{(c)}\le w+1\}|2^{-w/\sqrt8}\\
&\le2e^{-c_IA}(A+1)^Q\sum_{w=1}^A2^{-w/\sqrt8}
\le8(A+1)^Qe^{-c_IA}\\
&\le2^{(2/\log c)(\log\log n)^2}e^{-c_IA}\le\frac1{8n}.
\end{align*}
To finish bounding $\Delta_1$, it is left to bound the first term of Equation(11), denoted by $\Delta_{1,1}$.
\begin{authorclaim}{Claim23}
$\Delta_{1,1}=2|\sum_{\vec k\in\mathbb N^Q,\ 0<\|\vec k\|_{(c)}\le A}\mathbb E[S_{\vec k}(\vec g(Y))]|\le1/(8n)$.
\end{authorclaim}
\begin{proof}
The proof goes by bounding the spectral norm of the function $S_{\vec k}(\vec g(x))$. As for every $\vec k\in\mathbb N^Q$ with $\|\vec k\|_{(c)}\ne0$, $\mathbb E[S_{\vec k}(\vec g(U))]=0$, the claim will follow by using Claim5, together with sub-additivity and sub-multiplicativity. First, note that:
\begin{authorclaim}{Claim24}
For every $i\in[m]$, $L_1(\widehat g_i)\le4$.
\end{authorclaim}
\begin{proof}
Consider the function $h_i=1-g_i$, so $L_1(\widetilde g_i)\le L_1(\widetilde h_i)+1$ and $\mathbb E[\widetilde h_i]=\mathbb E[h_i]=1-\mathbb E[g_i]\ge1/2$. Now,
$L_1(\widehat h_i)\le(1/\mathbb E[\widetilde h_i])L_1(\widetilde h_i)+1\le2L_1(\widetilde h_i)+1$. Recalling that $\widetilde h_i(x)=\mathbb E[h_i(x+z\wedge U)]$, we get $L_1(\widetilde h_i)\le1$ as every shift of $h_i$ is a negated conjunction of literals. Thus, $L_1(\widehat h_i)\le3$ and $L_1(\widehat g_i)\le4$.
\end{proof}
Then, for every such $\vec k\in\mathbb N^Q$,
\begin{align*}
L_1(S_{\vec k}(\vec g))
&\le\sum_{\substack{I\subseteq\mathcal C\setminus X_T\\K(I)=\vec k}}\prod_{i\in I}L_1(\widehat g_i)
\le\sum_{\substack{I\subseteq\mathcal C\setminus X_T\\K(I)=\vec k}}4^{|I|}\\
&=\sum_{\substack{I\subseteq\mathcal C\setminus X_T\\K(I)=\vec k}}\prod_{q\in[Q]}4^{\vec k[q]}
=\prod_{q\in[Q]}\binom{\vec k[q]}{m_q}4^{\vec k[q]}.
\end{align*}
Recall that Claim18 tells us that $m_q\le3^{c^q}$, so
\begin{equation}
L_1(S_{\vec k}(\vec g))\le\prod_{q\in[Q]}3^{c^q(1+\log_34)\vec k[q]}\le12^{\|\vec k\|_{(c)}}.\tag{12}
\end{equation}
Finally,
\[
L_1\left(\sum_{\substack{\vec k\in\mathbb N^Q\\0<\|\vec k\|_{(c)}\le A}}\mathbb E[S_{\vec k}(\vec g(Y))]\right)
\le(A+1)^Q\cdot12^A\le2^{6A}\le n^{6\bar C},
\]
and the claim follows by observing that $\delta\le(1/(32n))n^{-6\bar C}$.
\end{proof}
Incorporating the above claim, we get that $\Delta_1=\Delta_{1,1}+\Delta_{1,2}\le1/(8n)+1/(8n)\le1/(4n)$, which readily gives $\Delta\le1/(4n)+2^{-A/1024}\le1/(2n)$. Plugging-it in Equation(9), we finally get
\[
|\mathbb E[\widetilde g(Y)]-\mathbb E[\widetilde g]|
\le\frac1{2n}+2e^{-c_IA}\le\frac1n
\]
and the desired result.
\end{proof}
